\documentclass[12pt, a4paper]{article}

% ── PAQUETES ─────────────────────────────────────────────────────────────────
\usepackage[spanish]{babel}
\usepackage[utf8]{inputenc}
\usepackage[T1]{fontenc}
\usepackage[margin=2.5cm]{geometry}
\usepackage{graphicx}
\usepackage{float}
\usepackage{hyperref}
\usepackage{booktabs}
\usepackage{array}
\usepackage{xcolor}
\usepackage{enumitem}
\usepackage{titlesec}
\usepackage{fancyhdr}
\usepackage{parskip}

% ── CONFIGURACIÓN HYPERREF ────────────────────────────────────────────────────
\hypersetup{
    colorlinks=true,
    linkcolor=blue!70!black,
    urlcolor=blue!70!black,
    pdftitle={Manual de Usuario - Sistema de Gestión de Notas},
    pdfauthor={Proyecto Notas v2}
}

% ── ENCABEZADO Y PIE DE PÁGINA ───────────────────────────────────────────────
\pagestyle{fancy}
\fancyhf{}
\rhead{Sistema de Gestión de Notas}
\lhead{\leftmark}
\rfoot{\thepage}

% ── COLORES DEL SISTEMA ───────────────────────────────────────────────────────
\definecolor{admincolor}{RGB}{37, 99, 235}    % azul
\definecolor{profcolor}{RGB}{5, 150, 105}     % verde
\definecolor{estudcolor}{RGB}{217, 119, 6}    % naranja
\definecolor{nocolor}{RGB}{220, 38, 38}       % rojo
\definecolor{yescolor}{RGB}{22, 163, 74}      % verde

% ── COMANDOS PERSONALIZADOS ───────────────────────────────────────────────────
\newcommand{\puede}{\textcolor{yescolor}{\textbf{Sí}}}
\newcommand{\nopuede}{\textcolor{nocolor}{\textbf{No}}}
\newcommand{\rol}[2]{\textcolor{#1}{\textbf{#2}}}

% ─────────────────────────────────────────────────────────────────────────────
\begin{document}

% ── PORTADA ──────────────────────────────────────────────────────────────────
\begin{titlepage}
    \centering
    \vspace*{2cm}

    % AQUÍ COLOCAR EL LOGO DE LA INSTITUCIÓN
    % Captura sugerida: logotipo oficial de la escuela, fondo blanco, formato PNG.
    % Sirve para identificar a qué institución pertenece el sistema.
    \includegraphics[width=0.35\textwidth]{logo_institucion.png}

    \vspace{1.5cm}
    {\LARGE \textbf{Sistema de Gestión de Notas}}\\[0.5cm]
    {\Large \textit{Proyecto Notas v2}}\\[0.3cm]
    \rule{\linewidth}{0.5pt}\\[0.8cm]
    {\Huge \textbf{Manual de Usuario}}\\[0.8cm]
    \rule{\linewidth}{0.5pt}

    \vspace{1.5cm}
    \begin{tabular}{ll}
        \textbf{Autor:}   & [Nombre del Autor] \\[0.3cm]
        \textbf{Versión:} & 1.0 \\[0.3cm]
        \textbf{Fecha:}   & \today \\[0.3cm]
    \end{tabular}

    \vfill
    {\small Documento de uso interno de la institución educativa}
\end{titlepage}

% ── ÍNDICE ───────────────────────────────────────────────────────────────────
\newpage
\tableofcontents
\newpage

% ─────────────────────────────────────────────────────────────────────────────
\section{Introducción}

El \textbf{Sistema de Gestión de Notas} es una plataforma web diseñada para instituciones educativas de nivel básico y bachillerato. Permite gestionar el ciclo académico completo: desde la creación de cursos y matrículas hasta el registro de calificaciones, cálculo de promedios y subida de evidencias en formato PDF.

El sistema opera enteramente desde el navegador web, sin necesidad de instalar programas adicionales en el equipo del usuario.

\subsection{Propósito del documento}

Este manual está dirigido a los tres tipos de usuarios del sistema: \rol{admincolor}{Administrador}, \rol{profcolor}{Profesor} y \rol{estudcolor}{Estudiante}. Describe de forma paso a paso cada funcionalidad disponible, así como las restricciones de acceso por rol.

% ─────────────────────────────────────────────────────────────────────────────
\section{Requisitos de Acceso}

Para utilizar el sistema no se requiere instalar ningún software adicional. Los requisitos mínimos son:

\begin{itemize}
    \item Navegador web actualizado: Google Chrome, Mozilla Firefox, Microsoft Edge o Safari (versión de los últimos 2 años)
    \item Conexión estable a internet
    \item Correo electrónico institucional registrado en el sistema
    \item Dispositivo: computador de escritorio, portátil o tablet (la pantalla de celular puede limitar la visualización de tablas de calificaciones)
\end{itemize}

% ─────────────────────────────────────────────────────────────────────────────
\section{Primeros Pasos}

\subsection{Inicio de sesión}

\begin{enumerate}
    \item Abre el navegador y dirígete a la dirección del sistema:\\
          \url{https://notas-v2-five.vercel.app}
    \item En la pantalla de inicio de sesión (Figura \ref{fig:login}), ingresa:
    \begin{itemize}
        \item \textbf{Correo electrónico:} tu email institucional (ejemplo: \texttt{admin@escuela.ec})
        \item \textbf{Contraseña:} la contraseña asignada por el administrador
    \end{itemize}
    \item Haz clic en el botón \textbf{``Ingresar''}.
    \item El sistema te redirigirá automáticamente al panel de control correspondiente a tu rol.
\end{enumerate}

\begin{figure}[H]
    \centering
    % AQUÍ COLOCAR: captura de pantalla de la pantalla de inicio de sesión.
    % Qué mostrar: el formulario con los campos "Correo electrónico" y "Contraseña",
    % el botón "Ingresar" y el enlace "¿Olvidaste tu contraseña?".
    % Resaltar con un recuadro rojo el campo de correo y con otro el de contraseña.
    % Por qué: es el primer contacto del usuario con el sistema; debe identificar
    % exactamente dónde ingresar sus datos.
    \includegraphics[width=0.75\textwidth]{login.png}
    \caption{Pantalla de inicio de sesión del sistema.}
    \label{fig:login}
\end{figure}

\subsection{Recuperación de contraseña}

Si el usuario olvidó su contraseña, puede recuperarla siguiendo estos pasos:

\begin{enumerate}
    \item En la pantalla de inicio de sesión, hacer clic en \textbf{``¿Olvidaste tu contraseña?''}
    \item Ingresar el correo electrónico institucional y hacer clic en \textbf{``Enviar''}
    \item Revisar la bandeja de entrada --- se recibirá un correo con un enlace de recuperación
    \item Hacer clic en el enlace del correo (tiene validez por tiempo limitado)
    \item Ingresar y confirmar la nueva contraseña
    \item Volver al inicio de sesión con las nuevas credenciales
\end{enumerate}

% ─────────────────────────────────────────────────────────────────────────────
\section{Descripción General de la Interfaz}

Al iniciar sesión correctamente, el sistema muestra el \textbf{panel de control} (dashboard). La interfaz se divide en tres zonas principales:

\begin{figure}[H]
    \centering
    % AQUÍ COLOCAR: captura del dashboard principal tras iniciar sesión como Administrador.
    % Qué mostrar: menú lateral izquierdo con todas las opciones, área central con
    % resumen de estadísticas (total de usuarios, cursos activos, etc.) y la
    % barra superior con el nombre del usuario y el botón de cerrar sesión.
    % Resaltar con flechas numeradas: 1) Menú lateral, 2) Área de contenido,
    % 3) Barra superior.
    % Por qué: el usuario debe orientarse en la interfaz antes de empezar a usarla.
    \includegraphics[width=\textwidth]{dashboard.png}
    \caption{Panel de control principal (vista del Administrador).}
    \label{fig:dashboard}
\end{figure}

\begin{enumerate}
    \item \textbf{Menú lateral (izquierda):} contiene los accesos a todos los módulos disponibles según el rol del usuario. Solo se muestran las opciones permitidas para cada rol.
    \item \textbf{Área de contenido (centro):} muestra la información del módulo seleccionado --- listas, formularios, tablas de calificaciones, etc.
    \item \textbf{Barra superior:} muestra el nombre del usuario activo y el botón para cerrar sesión.
\end{enumerate}

% ─────────────────────────────────────────────────────────────────────────────
\section{Roles del Sistema}

El sistema tiene tres roles con accesos diferenciados. Cada usuario tiene asignado un rol al momento de ser registrado por el administrador.

\begin{table}[H]
    \centering
    \begin{tabular}{>{\bfseries}p{3cm} p{10cm}}
        \toprule
        Rol & Descripción \\
        \midrule
        \rol{admincolor}{Administrador} & Gestiona la estructura académica: usuarios, cursos, materias, matrículas y docencias. Tiene el mayor nivel de acceso. \\[0.4cm]
        \rol{profcolor}{Profesor}       & Gestiona el contenido académico de sus materias: parciales, actividades, calificaciones y evidencias. \\[0.4cm]
        \rol{estudcolor}{Estudiante}    & Consulta sus propias calificaciones y promedios, y sube evidencias de sus actividades. \\
        \bottomrule
    \end{tabular}
    \caption{Roles del sistema y sus responsabilidades.}
    \label{tab:roles}
\end{table}

% ─────────────────────────────────────────────────────────────────────────────
\section{Guía de Funcionalidades: Administrador}
\label{sec:admin}

\subsection{Permisos del Administrador}

\begin{table}[H]
    \centering
    \begin{tabular}{p{5cm} c p{7cm}}
        \toprule
        \textbf{Acción} & \textbf{¿Permitido?} & \textbf{Detalle} \\
        \midrule
        Crear / editar / eliminar usuarios     & \puede    & Gestión completa de cuentas \\
        Asignar roles a usuarios               & \puede    & Administrador, Profesor, Estudiante \\
        Gestionar años lectivos                & \puede    & Crear, editar estado, eliminar \\
        Gestionar cursos                       & \puede    & Crear, editar, eliminar, asignar materias \\
        Gestionar materias                     & \puede    & Crear, editar, eliminar \\
        Matricular / desmatricular estudiantes & \puede    & Por curso \\
        Asignar / remover docencias            & \puede    & Profesor--Materia--Curso \\
        Ejecutar job de retención de evidencias & \puede   & Solo en años lectivos finalizados \\
        Crear actividades académicas           & \nopuede  & Exclusivo del Profesor \\
        Registrar calificaciones               & \nopuede  & Exclusivo del Profesor \\
        Subir evidencias en PDF                & \nopuede  & Exclusivo del Estudiante \\
        \bottomrule
    \end{tabular}
    \caption{Permisos del rol Administrador.}
\end{table}

\subsection{Gestionar usuarios}

\begin{enumerate}
    \item En el menú lateral seleccionar \textbf{``Usuarios''}.
    \item La pantalla mostrará la lista de todos los usuarios registrados (Figura \ref{fig:usuarios}).
    \item Para \textbf{crear un usuario}: hacer clic en \textbf{``Nuevo usuario''} y completar el formulario:
    \begin{itemize}
        \item Cédula (ID único de 10 dígitos)
        \item Nombre completo
        \item Correo electrónico institucional
        \item Contraseña temporal
        \item Rol: Administrador, Profesor o Estudiante
        \item Estado inicial: Activo
    \end{itemize}
    \item Para \textbf{editar} un usuario: hacer clic en el ícono de lápiz de la fila correspondiente.
    \item Para \textbf{cambiar el estado} (bloquear, inactivar): editar el usuario y modificar el campo de estado.
    \item Para \textbf{eliminar} un usuario: hacer clic en el ícono de papelera y confirmar la acción.
\end{enumerate}

\begin{figure}[H]
    \centering
    % AQUÍ COLOCAR: captura de la pantalla de listado de usuarios.
    % Qué mostrar: tabla con columnas Cédula, Nombre, Correo, Rol, Estado, y
    % columna de acciones con íconos de lápiz (editar) y papelera (eliminar).
    % Resaltar con un recuadro el botón "Nuevo usuario" en la esquina superior derecha.
    % Por qué: es la pantalla principal de gestión; el administrador debe identificar
    % rápidamente dónde están los controles de acción.
    \includegraphics[width=\textwidth]{usuarios_lista.png}
    \caption{Listado de usuarios con opciones de edición y eliminación.}
    \label{fig:usuarios}
\end{figure}

\noindent\textbf{Importante:} Eliminar un usuario es irreversible. Si solo se desea suspender el acceso temporalmente, cambiar el estado a \textit{Inactivo} o \textit{Bloqueado} en lugar de eliminar.

\subsection{Gestionar matrículas}

\begin{enumerate}
    \item Ir a \textbf{``Matrículas''} en el menú lateral.
    \item Seleccionar el curso.
    \item Hacer clic en \textbf{``Matricular estudiante''} y buscar al estudiante por nombre o cédula.
    \item Para desmatricular: seleccionar al estudiante en la lista y hacer clic en \textbf{``Desmatricular''}.
\end{enumerate}

\subsection{Asignar docencias}

\begin{enumerate}
    \item Ir a \textbf{``Docencias''} en el menú lateral.
    \item Hacer clic en \textbf{``Asignar docencia''}.
    \item Seleccionar el \textbf{profesor}, el \textbf{curso} y la \textbf{materia} que impartirá.
    \item Guardar. El profesor ya podrá ver ese curso-materia en su panel.
\end{enumerate}

% ─────────────────────────────────────────────────────────────────────────────
\section{Guía de Funcionalidades: Profesor}
\label{sec:profesor}

\subsection{Permisos del Profesor}

\begin{table}[H]
    \centering
    \begin{tabular}{p{5.5cm} c p{6.5cm}}
        \toprule
        \textbf{Acción} & \textbf{¿Permitido?} & \textbf{Detalle} \\
        \midrule
        Crear / eliminar parciales             & \puede    & Solo de sus materias asignadas \\
        Crear / editar / eliminar actividades  & \puede    & Tareas, exámenes, proyectos, lecciones \\
        Registrar calificaciones individuales  & \puede    & Con nota y comentario \\
        Calificar todo el curso a la vez       & \puede    & Función de calificación masiva \\
        Editar calificaciones ya registradas   & \puede    & \\
        Descargar evidencias de estudiantes    & \puede    & Solo de sus actividades \\
        Ver promedios y ranking del curso      & \puede    & \\
        Gestionar usuarios o cursos            & \nopuede  & Exclusivo del Administrador \\
        Matricular o desmatricular estudiantes & \nopuede  & Exclusivo del Administrador \\
        Subir evidencias en PDF                & \nopuede  & Exclusivo del Estudiante \\
        \bottomrule
    \end{tabular}
    \caption{Permisos del rol Profesor.}
\end{table}

\subsection{Crear actividades}

\begin{enumerate}
    \item En el menú seleccionar \textbf{``Actividades''}.
    \item Seleccionar el parcial al que pertenece la actividad.
    \item Hacer clic en \textbf{``Nueva actividad''} (Figura \ref{fig:nueva_actividad}).
    \item Completar el formulario:
    \begin{itemize}
        \item \textbf{Tipo:} Tarea / Examen / Proyecto / Lección
        \item \textbf{Título} de la actividad
        \item \textbf{Descripción} con las instrucciones
        \item \textbf{Fecha de inicio} y \textbf{fecha límite} de entrega
        \item \textbf{Valor máximo} en puntos
    \end{itemize}
    \item Guardar. Los estudiantes del curso recibirán una notificación automática.
\end{enumerate}

\begin{figure}[H]
    \centering
    % AQUÍ COLOCAR: captura del formulario de creación de una nueva actividad.
    % Qué mostrar: todos los campos del formulario visibles (Tipo, Título, Descripción,
    % Fechas, Valor máximo) y el botón "Guardar" al final.
    % Resaltar con flechas: el selector de Tipo (desplegable) y el campo de fechas,
    % ya que son los más propensos a errores de usuario.
    % Por qué: el profesor necesita ver exactamente qué datos debe ingresar para
    % que la actividad aparezca correctamente en el panel del estudiante.
    \includegraphics[width=0.8\textwidth]{nueva_actividad.png}
    \caption{Formulario de creación de nueva actividad.}
    \label{fig:nueva_actividad}
\end{figure}

\subsection{Registrar calificaciones}

\textbf{Calificación individual:}
\begin{enumerate}
    \item Ir a \textbf{``Calificaciones''} y seleccionar la actividad.
    \item Buscar al estudiante en la lista y hacer clic en \textbf{``Calificar''}.
    \item Ingresar la nota numérica y un comentario opcional.
    \item Guardar.
\end{enumerate}

\textbf{Calificación masiva (todo el curso):}
\begin{enumerate}
    \item Ir a \textbf{``Calificaciones''} y seleccionar la actividad.
    \item Hacer clic en \textbf{``Calificar curso completo''} (Figura \ref{fig:calificaciones_bulk}).
    \item Aparecerá la lista completa de estudiantes del curso. Ingresar la nota de cada uno.
    \item Hacer clic en \textbf{``Guardar todas''}.
\end{enumerate}

\begin{figure}[H]
    \centering
    % AQUÍ COLOCAR: captura de la pantalla de calificación masiva.
    % Qué mostrar: tabla con la columna de nombres de estudiantes y la columna de
    % entrada de nota para cada uno, más el botón "Guardar todas" al final.
    % Resaltar: la columna de notas con un recuadro azul para indicar que es editable.
    % Por qué: la calificación masiva es la función más usada y debe quedar clara
    % para que el profesor no calificar uno por uno innecesariamente.
    \includegraphics[width=\textwidth]{calificaciones_bulk.png}
    \caption{Pantalla de calificación masiva del curso completo.}
    \label{fig:calificaciones_bulk}
\end{figure}

% ─────────────────────────────────────────────────────────────────────────────
\section{Guía de Funcionalidades: Estudiante}
\label{sec:estudiante}

\subsection{Permisos del Estudiante}

\begin{table}[H]
    \centering
    \begin{tabular}{p{5.5cm} c p{6.5cm}}
        \toprule
        \textbf{Acción} & \textbf{¿Permitido?} & \textbf{Detalle} \\
        \midrule
        Ver sus propias calificaciones         & \puede    & Por actividad, parcial y materia \\
        Ver su promedio por materia y general  & \puede    & \\
        Ver su posición en el ranking del curso & \puede   & \\
        Subir evidencia en PDF                 & \puede    & Máx. 10 MB, un PDF por actividad \\
        Ver y eliminar sus propias evidencias  & \puede    & \\
        Ver y gestionar notificaciones         & \puede    & \\
        Ver calificaciones de otros estudiantes & \nopuede & \\
        Crear o editar actividades             & \nopuede  & Exclusivo del Profesor \\
        Registrar calificaciones               & \nopuede  & Exclusivo del Profesor \\
        Gestionar usuarios, cursos o materias  & \nopuede  & Exclusivo del Administrador \\
        \bottomrule
    \end{tabular}
    \caption{Permisos del rol Estudiante.}
\end{table}

\subsection{Ver mis calificaciones}

\begin{enumerate}
    \item Al iniciar sesión, el dashboard muestra un resumen de las notas recientes.
    \item Ir a \textbf{``Mis calificaciones''} en el menú lateral (Figura \ref{fig:mis_calificaciones}).
    \item Seleccionar el año lectivo y la materia para filtrar los resultados.
    \item La lista mostrará cada actividad con la nota obtenida y el comentario del profesor.
\end{enumerate}

\begin{figure}[H]
    \centering
    % AQUÍ COLOCAR: captura de la pantalla "Mis calificaciones" del estudiante.
    % Qué mostrar: tabla con columnas Actividad, Tipo, Fecha de entrega, Nota,
    % Comentario del profesor. Mostrar al menos 3 filas de ejemplo.
    % Resaltar: la columna de Nota con un recuadro verde para destacar la información
    % más buscada por el estudiante.
    % Por qué: es la pantalla más consultada por los estudiantes; debe verse de
    % forma clara y sin confusiones.
    \includegraphics[width=\textwidth]{mis_calificaciones.png}
    \caption{Pantalla de calificaciones del estudiante.}
    \label{fig:mis_calificaciones}
\end{figure}

\subsection{Subir una evidencia}

\begin{enumerate}
    \item Ir a \textbf{``Mis actividades''} en el menú lateral.
    \item Localizar la actividad para la que se desea subir la evidencia.
    \item Hacer clic en \textbf{``Subir evidencia''} (Figura \ref{fig:subir_evidencia}).
    \item Seleccionar el archivo PDF desde el dispositivo.
    \item Confirmar la subida.
\end{enumerate}

\begin{figure}[H]
    \centering
    % AQUÍ COLOCAR: captura del cuadro de diálogo o pantalla de subida de evidencia.
    % Qué mostrar: el botón de selección de archivo, el nombre del archivo seleccionado
    % y el botón de confirmación "Subir".
    % Resaltar: la indicación de "Solo PDF, máx. 10 MB" con una flecha y recuadro rojo.
    % Por qué: los estudiantes frecuentemente intentan subir archivos en formato
    % incorrecto; esta imagen les ayuda a ver la restricción antes de intentarlo.
    \includegraphics[width=0.7\textwidth]{subir_evidencia.png}
    \caption{Pantalla de subida de evidencia en PDF.}
    \label{fig:subir_evidencia}
\end{figure}

\noindent\textbf{Requisitos del archivo:}
\begin{itemize}
    \item Formato: \textbf{PDF únicamente}
    \item Tamaño máximo: \textbf{10 MB}
    \item Solo se permite \textbf{una evidencia por actividad}
\end{itemize}

\subsection{Eliminar y reemplazar una evidencia}

Si se subió un archivo incorrecto:

\begin{enumerate}
    \item Ir a \textbf{``Mis evidencias''} en el menú lateral.
    \item Localizar la evidencia que se desea eliminar.
    \item Hacer clic en el ícono de papelera y confirmar la eliminación.
    \item Una vez eliminada, volver a \textbf{``Mis actividades''} y subir el archivo correcto.
\end{enumerate}

% ─────────────────────────────────────────────────────────────────────────────
\section{Preguntas Frecuentes (FAQ)}

\subsection*{¿Por qué no puedo iniciar sesión si mi contraseña es correcta?}
La cuenta podría estar en estado \textit{Inactivo} o \textit{Bloqueado}. Contactar al administrador del sistema para que reactive el acceso.

\subsection*{¿Por qué no veo ningún curso ni materia en mi panel?}
\begin{itemize}
    \item \textbf{Profesor:} el administrador aún no ha asignado una docencia (Profesor--Materia--Curso).
    \item \textbf{Estudiante:} el administrador aún no ha realizado la matrícula en ningún curso.
\end{itemize}

\subsection*{¿Puedo subir evidencias en formato Word o imagen?}
No. El sistema acepta únicamente archivos en formato \textbf{PDF}. Se debe convertir el documento antes de subirlo.

\subsection*{¿Qué hago si subí la evidencia equivocada?}
Eliminarla desde \textbf{``Mis evidencias''} y subir el archivo correcto. Hacerlo antes de la fecha límite de entrega de la actividad.

\subsection*{¿Por qué el promedio no se actualiza después de que el profesor me calificó?}
El promedio se recalcula manualmente por el profesor después de registrar todas las calificaciones del parcial. Si persiste el problema, comunicarse con el profesor de la materia.

\subsection*{¿Puedo cambiar mi correo electrónico?}
No directamente. El correo electrónico solo puede ser modificado por el administrador del sistema.

\subsection*{¿Qué significa el estado ``Bloqueado'' en mi cuenta?}
El administrador ha restringido temporalmente el acceso. No será posible iniciar sesión hasta que el administrador cambie el estado a \textit{Activo}.

% ─────────────────────────────────────────────────────────────────────────────
\vfill
\begin{center}
    \rule{0.5\linewidth}{0.4pt}\\[0.3cm]
    {\small Para soporte técnico, contactar al administrador del sistema de la institución.}\\[0.2cm]
    {\small Sistema de Gestión de Notas --- Versión 1.0 --- \today}
\end{center}

\end{document}
